\documentclass{article}
\usepackage{amsmath, amssymb, geometry}
\geometry{a4paper, margin=1in}

\title{Homework 4.7 Variable-Coefficient Equations}
\author{Tashfeen Omran}
\date{September 2026}

\begin{document}
\maketitle

\section{Part 1: Comprehensive Introduction, Context, and Prerequisites}

\subsection{Core Concepts}
In previous chapters, we focused on linear second-order differential equations with \textit{constant} coefficients, taking the form $ay'' + by' + cy = f(t)$. However, many physical systems have properties that change over time or space, leading to equations where the coefficients are functions of the independent variable $t$. The standard form for a linear second-order variable-coefficient equation is:
\[ y'' + p(t)y' + q(t)y = g(t) \]
A special and highly important class of variable-coefficient equations is the \textbf{Cauchy-Euler equation} (or equidimensional equation), which takes the form:
\[ at^2 y'' + bty' + cy = f(t) \]
where $a$, $b$, and $c$ are constants. The defining characteristic of this equation is that the power of the independent variable $t$ matches the order of the derivative it multiplies.

\subsection{Intuition and Derivation}
For constant-coefficient equations, we guessed solutions of the form $y = e^{rt}$ because the derivative of an exponential is a multiple of itself, allowing the terms to cancel out. For Cauchy-Euler equations, this guess fails because the $t$ coefficients would remain. 
Instead, we exploit the equidimensional property by guessing a polynomial solution: $y = t^r$. 
Observe what happens when we take derivatives:
\[ y = t^r, \quad y' = r t^{r-1}, \quad y'' = r(r-1)t^{r-2} \]
Substituting these into the homogeneous Cauchy-Euler equation $at^2 y'' + bty' + cy = 0$ yields:
\[ a t^2 (r(r-1)t^{r-2}) + b t (r t^{r-1}) + c (t^r) = 0 \]
\[ t^r [a r(r-1) + br + c] = 0 \]
Assuming $t \neq 0$, we can divide by $t^r$ to obtain the \textbf{characteristic equation}:
\[ ar^2 + (b-a)r + c = 0 \]
This transforms the differential equation into a purely algebraic problem, exactly as we did for constant-coefficient equations!

\subsection{Historical Context and Motivation}
The study of variable-coefficient equations was heavily pioneered by Leonhard Euler and Augustin-Louis Cauchy in the 18th and 19th centuries. Mathematicians and physicists were attempting to solve Laplace's equation and the wave equation in cylindrical and spherical coordinate systems. When separating variables in these geometries, the resulting ordinary differential equations naturally acquire variable coefficients (specifically, terms like $r^2 R'' + r R'$). The Cauchy-Euler equation was one of the first variable-coefficient equations to be fully solved, serving as a stepping stone to more complex functions like Bessel functions and Legendre polynomials.

\subsection{Key Formulas}
\begin{itemize}
    \item \textbf{Standard Form:} $y'' + p(t)y' + q(t)y = g(t)$
    \item \textbf{Cauchy-Euler Characteristic Equation:} $ar^2 + (b-a)r + c = 0$
    \item \textbf{Cauchy-Euler Solutions ($t > 0$):}
    \begin{itemize}
        \item Distinct Real Roots ($r_1 \neq r_2$): $y(t) = c_1 t^{r_1} + c_2 t^{r_2}$
        \item Repeated Real Root ($r$): $y(t) = c_1 t^r + c_2 t^r \ln t$
        \item Complex Roots ($\alpha \pm i\beta$): $y(t) = t^\alpha (c_1 \cos(\beta \ln t) + c_2 \sin(\beta \ln t))$
    \end{itemize}
    \item \textbf{Reduction of Order:} Given one solution $y_1(t)$, a second linearly independent solution is:
    \[ y_2(t) = y_1(t) \int \frac{e^{-\int p(t) dt}}{y_1(t)^2} dt \]
    \item \textbf{Variation of Parameters:} For $y'' + p(t)y' + q(t)y = g(t)$, the particular solution is $y_p = v_1 y_1 + v_2 y_2$, where:
    \[ v_1 = \int \frac{-g(t)y_2(t)}{W} dt, \quad v_2 = \int \frac{g(t)y_1(t)}{W} dt \]
    \item \textbf{Abel's Identity:} $W[y_1, y_2](t) = C \exp\left[ -\int p(t) dt \right]$
\end{itemize}

\subsection{Prerequisites}
To succeed in this chapter, you must be proficient in:
1. Solving quadratic equations (including complex roots).
2. Integration techniques (u-substitution, integration by parts, trigonometric integrals).
3. The theory of linear independence and the Wronskian.
4. Euler's formula for complex exponentials: $e^{i\theta} = \cos\theta + i\sin\theta$.

---

\section{Part 2: Detailed Homework Solutions}
\textit{Note: A strategic selection of problems has been made to cover every distinct problem archetype, ensuring full conceptual mastery without redundant repetition.}

\subsection{Problem 2 (Existence and Uniqueness)}
\textbf{Question:} Determine the largest intervals for which Theorem 5 ensures the existence and uniqueness of a solution to $t(t-3)y'' + 2ty' - y = t^2$.

\textbf{Solution:}
First, we must put the equation into standard form $y'' + p(t)y' + q(t)y = g(t)$ by dividing by the leading coefficient $t(t-3)$:
\[ y'' + \frac{2}{t-3}y' - \frac{1}{t(t-3)}y = \frac{t}{t-3} \]
We identify the coefficient functions:
\[ p(t) = \frac{2}{t-3}, \quad q(t) = -\frac{1}{t(t-3)}, \quad g(t) = \frac{t}{t-3} \]
Next, we find the intervals where all three functions are simultaneously continuous.
The function $p(t)$ is discontinuous at $t = 3$.
The function $q(t)$ is discontinuous at $t = 0$ and $t = 3$.
The function $g(t)$ is discontinuous at $t = 3$.
The points of discontinuity for the system are $t = 0$ and $t = 3$. These points divide the real number line into three open intervals of continuity:
\[ (-\infty, 0), \quad (0, 3), \quad (3, \infty) \]
\textbf{Final Answer:} Theorem 5 ensures existence and uniqueness on any of the intervals $(-\infty, 0)$, $(0, 3)$, or $(3, \infty)$, depending on which interval contains the initial point $t_0$.

\subsection{Problem 9 (Cauchy-Euler, Distinct Real Roots)}
\textbf{Question:} Find a general solution to $t^2 y'' + 7ty' - 7y = 0$ for $t > 0$.

\textbf{Solution:}
This is a homogeneous Cauchy-Euler equation with $a=1, b=7, c=-7$.
We assume a solution of the form $y = t^r$. The characteristic equation is $ar^2 + (b-a)r + c = 0$:
\begin{align*}
(1)r^2 + (7 - 1)r - 7 &= 0 \\
r^2 + 6r - 7 &= 0
\end{align*}
Factoring the quadratic equation:
\[ (r+7)(r-1) = 0 \]
The roots are $r_1 = 1$ and $r_2 = -7$. Since the roots are real and distinct, the linearly independent solutions are $y_1 = t^1$ and $y_2 = t^{-7}$.
\textbf{Final Answer:} The general solution is $y(t) = c_1 t + c_2 t^{-7}$.

\subsection{Problem 14 (Cauchy-Euler, Repeated Real Roots)}
\textbf{Question:} Find a general solution to $t^2 y'' - 3ty' + 4y = 0$ for $t > 0$.

\textbf{Solution:}
This is a Cauchy-Euler equation with $a=1, b=-3, c=4$.
The characteristic equation is:
\begin{align*}
r^2 + (-3 - 1)r + 4 &= 0 \\
r^2 - 4r + 4 &= 0 \\
(r-2)^2 &= 0
\end{align*}
We have a repeated real root $r = 2$. The first solution is $y_1 = t^2$. For a repeated root in a Cauchy-Euler equation, the second linearly independent solution is obtained by multiplying the first by $\ln t$. Thus, $y_2 = t^2 \ln t$.
\textbf{Final Answer:} The general solution is $y(t) = c_1 t^2 + c_2 t^2 \ln t$.

\subsection{Problem 17 (Cauchy-Euler, Complex Roots)}
\textbf{Question:} Find a general solution to $t^2 y'' + 9ty' + 17y = 0$ for $t > 0$.

\textbf{Solution:}
Here, $a=1, b=9, c=17$. The characteristic equation is:
\begin{align*}
r^2 + (9 - 1)r + 17 &= 0 \\
r^2 + 8r + 17 &= 0
\end{align*}
We use the quadratic formula to find the roots:
\[ r = \frac{-8 \pm \sqrt{64 - 4(1)(17)}}{2} = \frac{-8 \pm \sqrt{64 - 68}}{2} = \frac{-8 \pm \sqrt{-4}}{2} = \frac{-8 \pm 2i}{2} = -4 \pm i \]
The roots are complex conjugates $\alpha \pm i\beta$ with $\alpha = -4$ and $\beta = 1$.
Using Euler's formula adapted for Cauchy-Euler equations, the solutions are $t^\alpha \cos(\beta \ln t)$ and $t^\alpha \sin(\beta \ln t)$.
\textbf{Final Answer:} The general solution is $y(t) = c_1 t^{-4} \cos(\ln t) + c_2 t^{-4} \sin(\ln t)$.

\subsection{Problem 15 (Cauchy-Euler for $t < 0$)}
\textbf{Question:} Find a general solution to $y'' - \frac{1}{t}y' + \frac{5}{t^2}y = 0$ for $t < 0$.

\textbf{Solution:}
First, multiply the entire equation by $t^2$ to put it into the standard Cauchy-Euler form:
\[ t^2 y'' - ty' + 5y = 0 \]
Here, $a=1, b=-1, c=5$. The characteristic equation is:
\begin{align*}
r^2 + (-1 - 1)r + 5 &= 0 \\
r^2 - 2r + 5 &= 0
\end{align*}
Using the quadratic formula:
\[ r = \frac{2 \pm \sqrt{4 - 20}}{2} = \frac{2 \pm 4i}{2} = 1 \pm 2i \]
We have complex roots with $\alpha = 1$ and $\beta = 2$.
Because the domain is $t < 0$, the function $\ln t$ is undefined. We must use the absolute value $|t| = -t$.
\textbf{Final Answer:} The general solution is $y(t) = c_1 (-t) \cos(2 \ln(-t)) + c_2 (-t) \sin(2 \ln(-t))$.

\subsection{Problem 19 (Initial Value Problem)}
\textbf{Question:} Solve the initial value problem $t^2 y'' - 4ty' + 4y = 0$, with $y(1) = -2$ and $y'(1) = -11$.

\textbf{Solution:}
The characteristic equation for this Cauchy-Euler equation ($a=1, b=-4, c=4$) is:
\begin{align*}
r^2 + (-4 - 1)r + 4 &= 0 \\
r^2 - 5r + 4 &= 0 \\
(r-1)(r-4) &= 0
\end{align*}
The roots are $r = 1$ and $r = 4$. The general solution and its derivative are:
\begin{align*}
y(t) &= c_1 t + c_2 t^4 \\
y'(t) &= c_1 + 4c_2 t^3
\end{align*}
Apply the initial conditions at $t=1$:
\begin{align*}
y(1) &= c_1(1) + c_2(1)^4 = c_1 + c_2 = -2 \\
y'(1) &= c_1 + 4c_2(1)^3 = c_1 + 4c_2 = -11
\end{align*}
We have a system of two equations:
1) $c_1 + c_2 = -2$
2) $c_1 + 4c_2 = -11$
Subtracting the first equation from the second gives:
\[ 3c_2 = -9 \implies c_2 = -3 \]
Substitute $c_2 = -3$ into the first equation:
\[ c_1 - 3 = -2 \implies c_1 = 1 \]
\textbf{Final Answer:} The unique solution to the IVP is $y(t) = t - 3t^4$.

\subsection{Problem 37 (Variation of Parameters)}
\textbf{Question:} Find a general solution to $t^2 z'' + tz' + 9z = -\tan(3 \ln t)$.

\textbf{Solution:}
\textbf{Step 1: Solve the homogeneous equation.}
$t^2 z'' + tz' + 9z = 0$. The characteristic equation is:
\[ r^2 + (1 - 1)r + 9 = 0 \implies r^2 + 9 = 0 \implies r = \pm 3i \]
The homogeneous solutions are $z_1 = \cos(3 \ln t)$ and $z_2 = \sin(3 \ln t)$.

\textbf{Step 2: Put the nonhomogeneous equation in standard form.}
Divide by $t^2$:
\[ z'' + \frac{1}{t}z' + \frac{9}{t^2}z = -\frac{\tan(3 \ln t)}{t^2} \]
Thus, $g(t) = -\frac{\tan(3 \ln t)}{t^2}$.

\textbf{Step 3: Compute the Wronskian.}
\begin{align*}
W(z_1, z_2) &= z_1 z_2' - z_1' z_2 \\
&= \cos(3 \ln t) \left( \frac{3}{t} \cos(3 \ln t) \right) - \left( -\frac{3}{t} \sin(3 \ln t) \right) \sin(3 \ln t) \\
&= \frac{3}{t} \cos^2(3 \ln t) + \frac{3}{t} \sin^2(3 \ln t) = \frac{3}{t} (\cos^2(3 \ln t) + \sin^2(3 \ln t)) = \frac{3}{t}
\end{align*}

\textbf{Step 4: Apply Variation of Parameters.}
\begin{align*}
v_1 &= \int \frac{-g(t)z_2(t)}{W} dt = \int \frac{-\left(-\frac{\tan(3 \ln t)}{t^2}\right) \sin(3 \ln t)}{3/t} dt \\
&= \int \frac{1}{3t} \tan(3 \ln t) \sin(3 \ln t) dt
\end{align*}
Let $u = 3 \ln t$, then $du = \frac{3}{t} dt$, which means $\frac{1}{3t} dt = \frac{1}{9} du$.
\begin{align*}
v_1 &= \frac{1}{9} \int \tan u \sin u \, du = \frac{1}{9} \int \frac{\sin^2 u}{\cos u} du = \frac{1}{9} \int \frac{1 - \cos^2 u}{\cos u} du \\
&= \frac{1}{9} \int (\sec u - \cos u) du = \frac{1}{9} (\ln|\sec u + \tan u| - \sin u) \\
&= \frac{1}{9} \ln|\sec(3 \ln t) + \tan(3 \ln t)| - \frac{1}{9} \sin(3 \ln t)
\end{align*}
Now for $v_2$:
\begin{align*}
v_2 &= \int \frac{g(t)z_1(t)}{W} dt = \int \frac{\left(-\frac{\tan(3 \ln t)}{t^2}\right) \cos(3 \ln t)}{3/t} dt \\
&= \int -\frac{1}{3t} \tan(3 \ln t) \cos(3 \ln t) dt = \int -\frac{1}{3t} \sin(3 \ln t) dt
\end{align*}
Using the same substitution $u = 3 \ln t$:
\[ v_2 = \frac{1}{9} \int -\sin u \, du = \frac{1}{9} \cos u = \frac{1}{9} \cos(3 \ln t) \]

\textbf{Step 5: Construct the particular solution.}
\begin{align*}
z_p &= v_1 z_1 + v_2 z_2 \\
&= \left[ \frac{1}{9} \ln|\sec(3 \ln t) + \tan(3 \ln t)| - \frac{1}{9} \sin(3 \ln t) \right] \cos(3 \ln t) + \left[ \frac{1}{9} \cos(3 \ln t) \right] \sin(3 \ln t) \\
&= \frac{1}{9} \cos(3 \ln t) \ln|\sec(3 \ln t) + \tan(3 \ln t)| - \frac{1}{9} \sin(3 \ln t) \cos(3 \ln t) + \frac{1}{9} \cos(3 \ln t) \sin(3 \ln t) \\
&= \frac{1}{9} \cos(3 \ln t) \ln|\sec(3 \ln t) + \tan(3 \ln t)|
\end{align*}
\textbf{Final Answer:} $z(t) = c_1 \cos(3 \ln t) + c_2 \sin(3 \ln t) + \frac{1}{9} \cos(3 \ln t) \ln|\sec(3 \ln t) + \tan(3 \ln t)|$.

\subsection{Problem 41 (Reduction of Order)}
\textbf{Question:} Given the differential equation $t^2 y'' - 2ty' - 4y = 0$ for $t > 0$ and a known solution $y_1(t) = t^{-1}$, find a second linearly independent solution using reduction of order.

\textbf{Solution:}
First, put the equation in standard form by dividing by $t^2$:
\[ y'' - \frac{2}{t}y' - \frac{4}{t^2}y = 0 \]
Identify $p(t) = -\frac{2}{t}$.
Calculate the integrating factor for the reduction of order formula:
\[ e^{-\int p(t) dt} = e^{-\int \left(-\frac{2}{t}\right) dt} = e^{2 \ln t} = t^2 \]
Apply the reduction of order formula:
\begin{align*}
y_2(t) &= y_1(t) \int \frac{e^{-\int p(t) dt}}{y_1(t)^2} dt \\
&= t^{-1} \int \frac{t^2}{(t^{-1})^2} dt \\
&= t^{-1} \int \frac{t^2}{t^{-2}} dt = t^{-1} \int t^4 dt \\
&= t^{-1} \left( \frac{t^5}{5} \right) = \frac{t^4}{5}
\end{align*}
Since any constant multiple of a solution is also a solution to a homogeneous linear equation, we can drop the $1/5$ to get the simplest form.
\textbf{Final Answer:} A second linearly independent solution is $y_2(t) = t^4$.

\subsection{Problem 32 (Abel's Formula)}
\textbf{Question:} By completing the following steps, prove that the Wronskian of any two solutions $y_1, y_2$ to the equation $y'' + p(t)y' + q(t)y = 0$ is given by Abel's formula: $W(t) = C \exp\left[ -\int p(t) dt \right]$.

\textbf{Solution:}
\textbf{(a) Show that the Wronskian $W$ satisfies the equation $W' + pW = 0$.}
By definition, $W = y_1 y_2' - y_1' y_2$.
Differentiating with respect to $t$ using the product rule:
\[ W' = (y_1' y_2' + y_1 y_2'') - (y_1'' y_2 + y_1' y_2') = y_1 y_2'' - y_1'' y_2 \]
Since $y_1$ and $y_2$ are solutions to the ODE, we can isolate their second derivatives:
\[ y_1'' = -p y_1' - q y_1 \]
\[ y_2'' = -p y_2' - q y_2 \]
Substitute these into the expression for $W'$:
\begin{align*}
W' &= y_1(-p y_2' - q y_2) - y_2(-p y_1' - q y_1) \\
&= -p y_1 y_2' - q y_1 y_2 + p y_1' y_2 + q y_1 y_2 \\
&= -p y_1 y_2' + p y_1' y_2 \\
&= -p(y_1 y_2' - y_1' y_2) \\
&= -p W
\end{align*}
Thus, $W' + pW = 0$.

\textbf{(b) Solve the separable equation in part (a).}
We have a first-order separable differential equation:
\[ \frac{dW}{dt} = -p(t)W \implies \frac{dW}{W} = -p(t) dt \]
Integrating both sides:
\[ \ln|W| = -\int p(t) dt + K \]
Exponentiating both sides:
\[ |W| = e^K \exp\left[ -\int p(t) dt \right] \implies W(t) = C \exp\left[ -\int p(t) dt \right] \]
where $C = \pm e^K$ is a constant determined by initial conditions.

\textbf{(c) How does Abel's formula clarify that the Wronskian is either identically zero or never zero?}
The exponential function $\exp(x)$ is strictly positive for all real numbers $x$. Therefore, the term $\exp\left[ -\int p(t) dt \right]$ can never equal zero. The only way the Wronskian $W(t)$ can be zero is if the constant $C$ is exactly $0$. If $C = 0$, then $W(t) = 0$ for all $t$ (identically zero). If $C \neq 0$, then $W(t) \neq 0$ for all $t$ (never zero).

---

\section{Part 3: In-Depth Analysis of Problems and Techniques}

\subsection{A) Problem Types and General Approach}
\begin{itemize}
    \item \textbf{Existence and Uniqueness Analysis (Problem 2):} The strategy is to force the equation into standard form, identify $p(t), q(t), g(t)$, and find the intersection of their domains of continuity. The theorem guarantees a unique solution on the specific continuous interval that contains the initial condition $t_0$.
    \item \textbf{Homogeneous Cauchy-Euler Equations (Problems 9, 14, 17, 15, 19):} The universal approach is to assume $y = t^r$, which transforms the differential equation into the algebraic characteristic equation $ar^2 + (b-a)r + c = 0$. The nature of the roots (distinct real, repeated real, or complex) dictates the template for the general solution.
    \item \textbf{Nonhomogeneous Cauchy-Euler Equations (Problem 37):} These require a two-phase approach. First, solve the associated homogeneous equation. Second, put the equation in standard form (crucial step!) and apply Variation of Parameters.
    \item \textbf{Reduction of Order (Problem 41):} When one solution is known, this technique constructs a second linearly independent solution. The approach relies on putting the equation in standard form to correctly identify $p(t)$ before evaluating the integral formula.
\end{itemize}

\subsection{B) Key Algebraic and Calculus Manipulations}
\begin{itemize}
    \item \textbf{The Equidimensional Property:} The core algebraic trick of this chapter is recognizing that $t^k \frac{d^k}{dt^k}(t^r)$ always results in a constant multiple of $t^r$. This perfectly balances the powers of $t$, allowing them to be factored out and eliminated.
    \item \textbf{Euler's Formula for Logarithms:} In Problem 17, we encountered $t^{-4 \pm i}$. By rewriting $t^i$ as $e^{i \ln t}$, we utilized Euler's formula ($e^{i\theta} = \cos\theta + i\sin\theta$) to convert complex exponentials into real-valued trigonometric functions of $\ln t$.
    \item \textbf{Trigonometric Simplification in Integration:} In Problem 37, the Variation of Parameters integral resulted in $\int \tan u \sin u \, du$. The necessary algebraic manipulation was rewriting this as $\frac{\sin^2 u}{\cos u}$, then using the Pythagorean identity to get $\frac{1 - \cos^2 u}{\cos u} = \sec u - \cos u$, which is easily integrable.
    \item \textbf{Domain Adaptation ($t < 0$):} In Problem 15, the domain was restricted to negative numbers. The critical manipulation was replacing $t$ with $-t$ (or $|t|$) inside the natural logarithm to prevent evaluating the logarithm of a negative number, ensuring the solution remained in the real domain.
\end{itemize}

---

\section{Part 4: ``Cheatsheet'' and Tips for Success}

\subsection{Summary of Important Formulas}
\begin{itemize}
    \item \textbf{Standard Form:} $y'' + p(t)y' + q(t)y = g(t)$
    \item \textbf{Cauchy-Euler Char. Eq.:} $ar^2 + (b-a)r + c = 0$
    \item \textbf{Cauchy-Euler Roots:}
    \begin{itemize}
        \item Real Distinct: $y = c_1 t^{r_1} + c_2 t^{r_2}$
        \item Real Repeated: $y = c_1 t^r + c_2 t^r \ln t$
        \item Complex ($\alpha \pm i\beta$): $y = t^\alpha (c_1 \cos(\beta \ln t) + c_2 \sin(\beta \ln t))$
    \end{itemize}
    \item \textbf{Reduction of Order:} $y_2 = y_1 \int \frac{e^{-\int p(t) dt}}{y_1^2} dt$
\end{itemize}

\subsection{Cheats, Tricks, and Shortcuts}
\begin{itemize}
    \item \textbf{The $t = e^x$ Cheat Code:} If you ever forget the Cauchy-Euler formulas during an exam, substitute $t = e^x$ (which means $x = \ln t$). This substitution magically transforms any Cauchy-Euler equation into a constant-coefficient equation, which you can solve using methods from previous chapters!
    \item \textbf{Wronskian Shortcut:} For a Cauchy-Euler equation, Abel's formula tells us that the Wronskian is proportional to $\exp(-\int \frac{b}{at} dt) = \exp(-\frac{b}{a} \ln t) = t^{-b/a}$. You can use this to check your manual Wronskian calculations.
\end{itemize}

\subsection{Common Pitfalls and Traps}
\begin{itemize}
    \item \textbf{The $(b-a)$ Trap:} The most common mistake in this chapter is writing the characteristic equation as $ar^2 + br + c = 0$. You \textbf{must} remember to subtract $a$ from $b$ for the middle term!
    \item \textbf{The Standard Form Mandate:} Before using Variation of Parameters or Reduction of Order, you absolutely must divide the entire equation by the leading coefficient to get $y''$ by itself. Forgetting this will result in the wrong $p(t)$ or $g(t)$ and ruin the entire problem.
    \item \textbf{Logarithm Domain Trap:} If a problem specifies $t < 0$, you must write $\ln(-t)$ in your final answer. Writing $\ln t$ will be marked incorrect because it is undefined for negative numbers.
\end{itemize}

---

\section{Part 5: Conceptual Synthesis and The ``Big Picture''}

\subsection{A) Thematic Connections}
The central theme of this chapter is \textbf{adaptation through transformation}. In previous chapters, we developed powerful tools (like the characteristic equation and Variation of Parameters) for constant-coefficient environments. Here, we face a variable environment. Instead of inventing entirely new math, we use the equidimensional property of polynomials to transform the variable-coefficient problem back into an algebraic one. Furthermore, the Superposition Principle remains intact; the fact that we can still build general solutions from linear combinations of $y_1$ and $y_2$ demonstrates that the core structure of linear vector spaces is preserved, even when the coefficients vary.

\subsection{B) Forward and Backward Links}
\textbf{Backward Links:} This topic is a direct, logical evolution of Chapter 4.2 (Constant-Coefficient Equations). It takes the algebraic characteristic equation method and adapts it for polynomial coefficients. It also heavily relies on the Variation of Parameters technique learned in Chapter 4.6. \\
\textbf{Forward Links:} Cauchy-Euler equations are the gateway to \textbf{Power Series Solutions} (Chapter 8). When solving advanced Partial Differential Equations (PDEs) like the Heat or Wave equation in cylindrical or spherical coordinates, the separation of variables frequently yields Cauchy-Euler equations or related forms like Bessel's equation. Understanding how to handle singular points (like $t=0$ in this chapter) is the foundational prerequisite for mastering the Method of Frobenius later in the curriculum.

---

\section{Part 6: Real-World Application and Modeling}

\subsection{A) Concrete Scenarios}
\begin{itemize}
    \item \textbf{Financial Mathematics (Option Pricing):} In quantitative finance, the Black-Scholes equation is used to price options. When pricing a \textit{perpetual} American put option (an option with no expiration date), the time derivative drops out. The resulting ordinary differential equation has coefficients that depend on the square of the stock price, perfectly matching the Cauchy-Euler form.
    \item \textbf{Fluid Mechanics (Couette Flow):} When analyzing the steady, axisymmetric flow of a viscous fluid between two rotating concentric cylinders, the Navier-Stokes equations in cylindrical coordinates simplify to a Cauchy-Euler equation for the angular velocity of the fluid as a function of radial distance.
    \item \textbf{Solid Mechanics (Pressure Vessels):} Calculating the radial and hoop stresses in a thick-walled cylindrical pipe subjected to internal and external pressures involves solving a Cauchy-Euler equation derived from the equilibrium equations in polar coordinates.
\end{itemize}

\subsection{B) Model Problem Setup: Perpetual American Put Option}
\textbf{Scenario:} A quantitative analyst at a hedge fund needs to determine the fair value of a perpetual American put option. \\
\textbf{Variables:}
\begin{itemize}
    \item $S$: The current price of the underlying stock (independent variable, $S > 0$).
    \item $V(S)$: The value of the perpetual put option (dependent variable).
    \item $\sigma$: The constant volatility of the stock.
    \item $r$: The constant risk-free interest rate.
\end{itemize}
\textbf{Model Formulation:}
The standard Black-Scholes partial differential equation is:
\[ \frac{\partial V}{\partial t} + \frac{1}{2}\sigma^2 S^2 \frac{\partial^2 V}{\partial S^2} + r S \frac{\partial V}{\partial S} - rV = 0 \]
Because the option is perpetual, its value does not depend on time $t$, meaning $\frac{\partial V}{\partial t} = 0$. The PDE reduces to a second-order linear ordinary differential equation:
\[ \frac{1}{2}\sigma^2 S^2 V'' + r S V' - rV = 0 \]
This is a homogeneous Cauchy-Euler equation! To find the option price, the analyst would solve this equation for $V(S)$ subject to financial boundary conditions (e.g., $V(S) \to 0$ as $S \to \infty$, and a smooth pasting condition at the optimal exercise boundary).

---

\section{Part 7: Common Variations and Untested Concepts}

\subsection{Untested Concept: Solving via the $t = e^x$ Substitution}
\textbf{Explanation:} A powerful technique not explicitly tested in the selected homework is transforming a Cauchy-Euler equation into a constant-coefficient equation using the substitution $t = e^x$ (or $x = \ln t$). This is incredibly useful for nonhomogeneous equations because it allows you to use the Method of Undetermined Coefficients instead of the often tedious Variation of Parameters.

\textbf{Worked Example:} Solve $t^2 y'' - 2ty' + 2y = t^3 \ln t$ for $t > 0$. \\
Let $t = e^x$, so $x = \ln t$. Using the chain rule, it can be shown that:
\[ t \frac{dy}{dt} = \frac{dy}{dx} \quad \text{and} \quad t^2 \frac{d^2y}{dt^2} = \frac{d^2y}{dx^2} - \frac{dy}{dx} \]
Substitute these into the original equation:
\begin{align*}
\left( \frac{d^2y}{dx^2} - \frac{dy}{dx} \right) - 2 \frac{dy}{dx} + 2y &= (e^x)^3 (x) \\
\frac{d^2y}{dx^2} - 3 \frac{dy}{dx} + 2y &= x e^{3x}
\end{align*}
This is now a constant-coefficient equation! 
The homogeneous roots are $r^2 - 3r + 2 = 0 \implies r=1, 2$, so $y_h(x) = c_1 e^x + c_2 e^{2x}$.
Using Undetermined Coefficients, we guess $y_p(x) = (Ax + B)e^{3x}$. 
Skipping the basic algebra, substituting this guess yields $A = 1/2$ and $B = -3/4$.
So, $y(x) = c_1 e^x + c_2 e^{2x} + \left( \frac{1}{2}x - \frac{3}{4} \right) e^{3x}$.
Finally, substitute back $e^x = t$ and $x = \ln t$ to get the solution in terms of $t$:
\[ y(t) = c_1 t + c_2 t^2 + \left( \frac{1}{2} \ln t - \frac{3}{4} \right) t^3 \]

---

\section{Part 8: Advanced Diagnostic Testing: ``Find the Flaw''}

\subsection{Problem 1}
\textbf{Task:} Solve $t^2 y'' + 3ty' + y = 0$ for $t > 0$. \\
\textbf{Flawed Solution:}
Assume $y = t^r$. The characteristic equation is $r^2 + 3r + 1 = 0$.
Using the quadratic formula, $r = \frac{-3 \pm \sqrt{9 - 4}}{2} = \frac{-3 \pm \sqrt{5}}{2}$.
The general solution is $y(t) = c_1 t^{(-3+\sqrt{5})/2} + c_2 t^{(-3-\sqrt{5})/2}$. \\
\textbf{Error Analysis:} The characteristic equation for a Cauchy-Euler equation $at^2 y'' + bty' + cy = 0$ is $ar^2 + (b-a)r + c = 0$, not $ar^2 + br + c = 0$. \\
\textbf{Correct Step:} The correct equation is $r^2 + (3-1)r + 1 = 0 \implies r^2 + 2r + 1 = 0 \implies (r+1)^2 = 0$. \\
\textbf{Correct Final Solution:} $y(t) = c_1 t^{-1} + c_2 t^{-1} \ln t$.

\subsection{Problem 2}
\textbf{Task:} Find the general solution to $t^2 y'' - 2y = t^3$ for $t > 0$. \\
\textbf{Flawed Solution:}
The homogeneous equation is $t^2 y'' - 2y = 0$. The characteristic equation is $r^2 - r - 2 = 0 \implies (r-2)(r+1) = 0$.
So $y_1 = t^2$ and $y_2 = t^{-1}$. The Wronskian is $W = t^2(-t^{-2}) - (2t)(t^{-1}) = -3$.
Using Variation of Parameters, $v_1 = \int \frac{-g(t)y_2(t)}{W} dt = \int \frac{-t^3(t^{-1})}{-3} dt = \int \frac{t^2}{3} dt = \frac{t^3}{9}$.
$v_2 = \int \frac{g(t)y_1(t)}{W} dt = \int \frac{t^3(t^2)}{-3} dt = \int -\frac{t^5}{3} dt = -\frac{t^6}{18}$.
$y_p = v_1 y_1 + v_2 y_2 = \frac{t^3}{9}(t^2) - \frac{t^6}{18}(t^{-1}) = \frac{t^5}{18}$. \\
\textbf{Error Analysis:} The function $g(t)$ in the Variation of Parameters formula must be taken from the standard form $y'' + p(t)y' + q(t)y = g(t)$, meaning we must divide the original equation by $t^2$ first. \\
\textbf{Correct Step:} Divide by $t^2$ to get $y'' - \frac{2}{t^2}y = t$, so $g(t) = t$. \\
\textbf{Correct Final Solution:} $v_1 = \int \frac{-t(t^{-1})}{-3} dt = \frac{t}{3}$. $v_2 = \int \frac{t(t^2)}{-3} dt = -\frac{t^4}{12}$. $y_p = \frac{t}{3}(t^2) - \frac{t^4}{12}(t^{-1}) = \frac{t^3}{4}$. General solution: $y(t) = c_1 t^2 + c_2 t^{-1} + \frac{t^3}{4}$.

\subsection{Problem 3}
\textbf{Task:} Use reduction of order to find a second linearly independent solution to $t^2 y'' - 3ty' + 4y = 0$ given $y_1 = t^2$. \\
\textbf{Flawed Solution:}
The reduction of order formula is $y_2 = y_1 \int \frac{e^{-\int p(t) dt}}{y_1^2} dt$.
Here, $p(t) = -3t$.
$\int p(t) dt = -\frac{3t^2}{2}$.
$e^{-\int p(t) dt} = e^{3t^2/2}$.
$y_2 = t^2 \int \frac{e^{3t^2/2}}{t^4} dt$. This integral cannot be evaluated in terms of elementary functions. \\
\textbf{Error Analysis:} The coefficient $p(t)$ must be identified after putting the differential equation into standard form by dividing by $t^2$. \\
\textbf{Correct Step:} Standard form is $y'' - \frac{3}{t}y' + \frac{4}{t^2}y = 0$, so $p(t) = -\frac{3}{t}$. \\
\textbf{Correct Final Solution:} $e^{-\int (-3/t) dt} = e^{3 \ln t} = t^3$. $y_2 = t^2 \int \frac{t^3}{(t^2)^2} dt = t^2 \int \frac{1}{t} dt = t^2 \ln t$.

\subsection{Problem 4}
\textbf{Task:} Solve $t^2 y'' - ty' + 5y = 0$ for $t > 0$. \\
\textbf{Flawed Solution:}
Assume $y = t^r$. The characteristic equation is $r^2 - 2r + 5 = 0$.
Roots are $r = \frac{2 \pm \sqrt{4 - 20}}{2} = 1 \pm 2i$.
The general solution is $y(t) = c_1 e^t \cos(2t) + c_2 e^t \sin(2t)$. \\
\textbf{Error Analysis:} The roots $r = \alpha \pm i\beta$ for a Cauchy-Euler equation yield solutions of the form $t^\alpha \cos(\beta \ln t)$ and $t^\alpha \sin(\beta \ln t)$, not exponentials. \\
\textbf{Correct Step:} Substitute $\alpha = 1$ and $\beta = 2$ into the correct Cauchy-Euler complex root template. \\
\textbf{Correct Final Solution:} $y(t) = c_1 t \cos(2 \ln t) + c_2 t \sin(2 \ln t)$.

\subsection{Problem 5}
\textbf{Task:} Solve $t^2 y'' + 5ty' + 4y = 0$ for $t < 0$. \\
\textbf{Flawed Solution:}
Assume $y = t^r$. The characteristic equation is $r^2 + 4r + 4 = 0 \implies (r+2)^2 = 0$.
The root is $r = -2$ with multiplicity 2.
The general solution is $y(t) = c_1 t^{-2} + c_2 t^{-2} \ln t$. \\
\textbf{Error Analysis:} For $t < 0$, the natural logarithm $\ln t$ is undefined; the correct second linearly independent solution must use $\ln(-t)$ to remain in the real domain. \\
\textbf{Correct Step:} Replace $\ln t$ with $\ln(-t)$ for the domain $t < 0$. \\
\textbf{Correct Final Solution:} $y(t) = c_1 t^{-2} + c_2 t^{-2} \ln(-t)$.

\end{document}